% ===================== PDF page 1 =====================
\printedpage{--}
% head: (none)
% no page number printed on this page

\begin{center}
\vspace*{5cm}

{\Large\textsf{THE DERIVATORS}}

\bigskip

\textbf{ALEXANDRE GROTHENDIECK}

\vspace{2cm}

\textbf{Chapter XVII}

\medskip

\textbf{Categories with fibrations and cofibrations}\\
\textbf{(Models (3))}
\end{center}

\clearpage

% ===================== PDF page 2 =====================
\printedpage{--}
% head: (none)
% no page number printed on this page

\vspace*{5cm}

This text was deciphered and transcribed in \LaTeXe{} by M. Künzer. It was edited by M. Künzer, J. Malgoire, and G. Maltsiniotis. The transcription is as faithful as possible to the manuscript. Only a few obvious minor corrections have been made. For the rare additions or comments by the editors, as well as for the original numbering of the manuscript pages, typewriter characters \texttt{[typewriter]} in brackets are used.

This edition is provisional. Remarks, comments and corrections are welcome. Send a message to:

\textbf{[e-mail address removed]}

\hspace*{\fill}G. Maltsiniotis\hspace*{2cm}

\clearpage

% ===================== PDF page 3 =====================
\printedpage{1}
% head: A. Grothendieck : Dérivateurs, Ch. XVII, Édition des Universités de Montpellier II et Paris VII
% foot: Version 9 février 2019

\msp{page 1}

\section*{1\quad Preliminaries\footnote{\textbf{N.B.} These preliminaries will be completely taken up again and superseded in section 2. Not to be read!}}

\textbf{Proposition 1.} \textit{Let $\mathcal{M}$ be a category stable under arbitrary $\varinjlim$, $C \subset \Fl(\mathcal{M})$ containing the isomorphisms, stable under cobase change and under composition} of transfinite length. \textit{Let, for every small ordered set $I$, $C_I \subset \Fl(\mathcal{M}(I^{\circ}))$ be the set of arrows $u_\bullet : X_\bullet \longrightarrow Y_\bullet$ such that $\forall\ i_0 \in I$, in the diagram} % ?? uncertain: superscript of I -- \circ or degree sign

\[
\begin{tikzcd}[column sep=large]
X'_{i_0} = \varinjlim\limits_{i<i_0} X_i \arrow[r] \arrow[d, "\alpha_{i_0} = \varinjlim\limits_{i<i_0} u_i"'] \arrow[dr, phantom, "\text{cocart.}"] & X_{i_0} \arrow[d, "\beta_{i_0}"'] \arrow[dd, dotted, shift left=3ex, "u_{i_0}"] & \\
Y'_{i_0} = \varinjlim\limits_{i<i_0} Y_i \arrow[r] & \tilde{Y}_{i_0} \arrow[d, "\tilde{u}_{i_0}"'] & \text{\footnotesize (defined as an amalgamated sum)} \\
& Y_{i_0}\ , &
\end{tikzcd}
\tag*{$(*)$}
\]

\textit{the arrow $\tilde{u}_{i_0}$ is in $C$\footnote{Condition stable under passage from $I$ to an \textit{open} ordered subset $I'$ of $I$.}. I say that if $I$ is} well ordered, \textit{then $u_\bullet \in C_I$ implies}
\begin{itemize}[leftmargin=2.5em]
\item[a)] \textit{that $\varinjlim u_i : \varinjlim X_i \longrightarrow \varinjlim Y_i$ is in $C$, and}
\item[b)] $\forall\ i \in I$, $(u_i : X_i \longrightarrow Y_i) \in C$\footnote{Say also that we have
\begin{itemize}[leftmargin=2em]
\item[b$'$)] \textit{If $i' \leq i$, then in the diagram similar to $(*)$,}
\end{itemize}
\[
\begin{tikzcd}[ampersand replacement=\&]
X'_{i'} \arrow[r] \arrow[d] \arrow[dr, phantom, "\text{cocart.}"] \& X'_i \arrow[d] \\
Y'_{i'} \arrow[r] \& \tilde{Y}'_i \arrow[d, "u_{i',i}"] \\
\& Y'_i\ ,
\end{tikzcd}
\]
\textit{we have} $u_{i',i} \in C$ (we recover $(*)$ by taking $i' = i_0$, $i = i_0 + 1$). Case b) is obtained by taking $i' = 0$, $i$ arbitrary, whence $(X'_i \longrightarrow Y'_i) \in C$. Applying it to $i+1$ instead of $i$, we find $(u_i : X_i \longrightarrow Y_i) \in C$.}.
\end{itemize}

. % a lone full stop printed on its own line

\textbf{Corollary.} \textit{Let $W \subset \mathcal{M}$ be a mildly saturated localizer, such that $TC \overset{\text{def}}{=} C \cap W$ is stable under cobase change, and under composition} of transfinite length. \textit{Let $u_\bullet : X_\bullet \longrightarrow Y_\bullet$ be in $C_I$. Then $u_\bullet \in W_{I^{\circ}}$ (i.e. $\forall\ i \in I$, we have $u_i \in W$) if and only if $u_\bullet \in (TC)_I$, i.e. if and only if for every $i \in I$, the arrow $\tilde{u}_i$ in $(*)$ is in $TC$, or what}
\msp{page 2 n'existe pas}
\printedpage{2}
% head: A. Grothendieck : Dérivateurs, Ch. XVII, Édition des Universités de Montpellier II et Paris VII
% foot: Version 9 février 2019
\msp{page 3}
\textit{amounts to the same thing, in $W$ (since by hypothesis it is already in $C$). If this condition is satisfied, then $\varinjlim u_i : \varinjlim X_i \longrightarrow \varinjlim Y_i$ is in $W$.}

% ===================== PDF page 4 (continued) =====================

\textsc{Proof} of the corollary. Suppose $u_\bullet \in (TC)_I$, then by the proposition, $\varinjlim\limits_{i\in I} u_i$ is in $TC$, or what amounts to the same thing, in $W$. On the other hand, applying this same result to the direct systems induced on $I_{<i_0}$, we find that in $(*)$, the arrow $\alpha_{i_0}$ is in $TC$, hence also $\beta_{i_0}$ by cobase change, hence also $u_{i_0} = \tilde{u}_{i_0}\beta_{i_0}$. So it remains to prove that conversely, if we assume the $u_{i_0} \in W$, then the $\tilde{u}_{i_0}$ are themselves also in $W$ (i.e. in $TC = C \cap W$), or what amounts to the same thing, that the $\beta_{i_0}$ are in $W$, or again in $TC$. But since $TC$ is stable under cobase change, it suffices for this to prove that $\alpha_{i_0} \in TC$, or again $\alpha_{i_0} \in C$. So we are reduced to proving that
\begin{itemize}[leftmargin=2.5em]
\item[c)] $u \in W_{I^{\circ}} \cap C_I \Longrightarrow (\,\varinjlim\limits_{i<i_0} u_i) \in W$ for every $i_0 \in I$,
\end{itemize}
so, changing notation, to proving \texttt{[incomplete sentence]}. We will prove a) b) c) together by transfinite induction on the ordinal number $\alpha$ such that $I \simeq I_\alpha = \{\text{set of ordinals } i < \alpha\}$.

\msp{page 4}

1$^\circ$) Case $\alpha = 0$, i.e. $I = \emptyset$: it is trivial.

2$^\circ$) If it's ok for $I_\alpha$, also for $I_{\alpha+1} = \{\beta \mid \beta \leq \alpha\}$. Here, $I_{\alpha+1}$ having a greatest element $\alpha$, a) is contained in b) (for $i = \alpha$, greatest element of $I_{\alpha+1}$.) To prove b), it remains to prove that $u_\alpha \in C$. We look at the diagram $(*)$ for $i_0 = \alpha$. By the induction hypothesis (form a), applied to $I_\alpha$, we have $\alpha_{i_0} \in C$, hence by cobase change $\beta_{i_0} \in C$, hence by composition $u_{i_0} = \tilde{u}_{i_0}\beta_{i_0} \in C$. To prove c), it remains to prove that $\varinjlim\limits_{i<\alpha} u_i \in W$. But by the induction hypothesis applied to $I_\alpha = \{i \mid i < \alpha\}$, we already know that $(u_i)_{i<\alpha} \in (TC)_{I_\alpha}$, hence by the proved part of the corollary or of the proposition, $\varinjlim\limits_{i<\alpha} u_i \in TC$.

3$^\circ$) If $\alpha$ is a limit ordinal, and if a) b) c) hold for $I_\beta$ with $\beta < \alpha$, they also hold for $I_\alpha$. This time, b) is contained in the induction
\msp{page 5}
hypothesis, we must prove a), i.e. that
\[
u_\alpha \overset{\text{def}}{=} \varinjlim_{i<\alpha} u_i \ :\ \underbrace{X_\alpha}_{\overset{\text{def}}{=} \varinjlim\limits_{i<\alpha} X_i} \longrightarrow \underbrace{Y_\alpha}_{\overset{\text{def}}{=} \varinjlim\limits_{i<\alpha} Y_i}
\]
is in $C$. As for c), it is already contained in the induction hypothesis. So it finally remains to prove
\[
u_\alpha \ \in \ C
\]
when we assume that the proposition is proved for the ordinals $\beta < \alpha$.

We are going to construct a third ordinal system on $I_\alpha$, by introducing for every $0 \leq i < \alpha$, i.e. $i \in I_\alpha$, $Z_i$ by the cocartesian square
\printedpage{3}
% head: A. Grothendieck : Dérivateurs, Ch. XVII, Édition des Universités de Montpellier II et Paris VII
% foot: Version 9 février 2019
\[
\begin{tikzcd}
X_i \arrow[r] \arrow[d] \arrow[dr, phantom, "\text{cocart.}"] & X_\alpha \arrow[d] \\
Y_i \arrow[r] & Z_i\ .
\end{tikzcd}
\]

The transition arrows between the $Z_i$ are obvious, and we will have $\varinjlim\limits_{i} Z_i \xrightarrow{\ \sim\ } Y_\alpha$. Thus, $u_\alpha : X_\alpha (= Z_{-1}) \longrightarrow Y_\alpha$ is decomposed into a transfinite composite of morphisms
\msp{page 6}
\[
Z_i \xrightarrow{v_{i+1}} Z_{i+1} \qquad (-1 \leq i < \alpha)
\tag*{$(*)$}
\]
and
\[
\varinjlim_{i<i_0} Z_i \xrightarrow{v_{i_0}} Z_{i_0} \qquad i_0 \text{ limit ordinal } < \alpha.
\tag*{$(**)$}
\]

(\textbf{N.B.} The case where $i_0$ is \textit{not} a limit ordinal is none other than $(*)$ for $i = \beta_0 - 1$.) Everything comes down to proving that these arrows are in $C$ (hence their transfinite composite too). % [sic] $\beta_0 - 1$ -> $i_0 - 1$, the case discussed is that of $i_0$

Case $(*)$.
\[
\begin{tikzcd}
X_i \arrow[r] \arrow[d, "u_i"'] \arrow[dr, phantom, "\text{cocart.}"] & X_{i+1} \arrow[r] \arrow[d, "\beta_{i+1}"'] \arrow[dr, phantom, "\text{cocart.}"] & X_\alpha \arrow[d] \\
Y_i \arrow[r] & \tilde{Y}_{i+1} \arrow[r] \arrow[d, "\tilde{u}_{i+1}"'] \arrow[dr, phantom, "\text{cocart.}"] & Z_i \arrow[d, "v_{i+1}"] \\
& Y_{i+1} \arrow[r] & Z_{i+1}\ .
\end{tikzcd}
\tag*{\footnotemark}
\]
\footnotetext{To give oneself confidence -- one could treat the general case right away.} % [sic] toute de suite -> tout de suite, French idiom

Since $\tilde{u}_{i+1} \in C$, $v_{i+1} \in C$ by cobase change, OK. We must however treat separately the case $i = 0$, i.e. the arrow $X_\alpha \longrightarrow Z_0$, which is included in
\[
\begin{tikzcd}
X_0 \arrow[r] \arrow[d, "u_0"'] \arrow[dr, phantom, "\text{cocart.}"] & X_\alpha = Z_{-1} \arrow[d, "v_0"] \\
Y_0 \arrow[r] & Z_0\ ,
\end{tikzcd}
\]
where $u_0 \in C$, hence $v_0 \in C$, OK.

Case $(**)$. If $i_0 < \alpha$ is a limit ordinal, let us inspect $(**)$. We have by passage to the limit
\[
Z'_{i_0} \overset{\text{def}}{=} \varinjlim_{i<i_0} Z_i \ \simeq\ Y'_{i_0} \amalg_{X'_{i_0}} X_\alpha
\]
(cf. diagram $(*)$ for the definition of $X'_{i_0}$, $Y'_{i_0}$), fitting into the following diagram grafted onto $(*)$